%=============================================================================!
% 18.6  STOCHASTIC GRADIENTS                                                  !
%=============================================================================!
\section{Stochastic gradients}
\label{sec:lo-sgd}

A poll does not ask every voter; a thousand chosen at random give the
share of a party to within a few per cent, at a tiny fraction of the
cost. The gradient of a loss is an average over the readings, and it too
can be estimated from a few of them. A learned potential is trained on
hundreds of thousands of configurations, each with hundreds of forces,
and an exact gradient over all of them for every step would make training
impossibly slow.

\subsection*{A gradient from a few readings}

Written as the mean over $n$ readings of a loss $\ell_i$ for each,
$L = \frac1n\sum_i\ell_i(\vb{w})$, the loss has the gradient $\grad L =
\frac1n\sum_i\grad\ell_i$, a mean of $n$ numbers. A \emph{mini-batch} of
$b$ readings chosen at random gives the estimate $\vb{g} =
\frac1b\sum_{i\in\text{batch}}\grad\ell_i$, whose average over the
possible batches is the true gradient. If the readings' gradients scatter
about their mean with the spread $\sigma_g$, each component of $\vb{g}$
scatters by
\begin{equation*}
   \frac{\sigma_g}{\sqrt b} ,
\end{equation*}
the error of a mean of $b$ independent samples (\cref{sec:cv-mean}).
\emph{Stochastic gradient descent} steps with $\vb{g}$ in place of
$\grad L$, one step per batch, and a pass through all the readings, an
\emph{epoch}, takes $n/b$ steps for the cost of one exact gradient.

\Cref{fig:lo-sgd} fits the force of the A--B bond at \SI{300}{\kelvin},
2001 readings from \cref{sec:lo-loss}, with a straight line in the
scaled stretch $x/\sigma_x$, its loss the mean squared miss. At the least
loss, where the exact gradient is zero, the gradient of the slope's
weight from batches of ten readings scatters by \num{0.180}, against the
$\sigma_g/\sqrt{10} = \num{0.185}$ that the readings' own spread gives
(\cref{fig:lo-sgd}b); its histogram is not yet Gaussian, since a mean of
ten skewed numbers is still skewed, but its width is as predicted.

\subsection*{A noise floor, and a rate that falls}

The noise has a cost. With the rate fixed at half the stability limit,
full-batch descent reaches the least loss in under ten epochs, while
batches of ten never settle: after 50 epochs their loss is still
\num{4.7e-3}\,(\si{\electronvolt\per\angstrom})$^2$ above the least, with
spikes far above it (\cref{fig:lo-sgd}a). Near the minimum the true
gradient vanishes but the estimate does not, and each step moves the
weights by about $\eta\sigma_g/\sqrt b$ in a random direction, so the
weights wander in a cloud whose size is set by the rate. A rate that
falls as training proceeds, here $\eta/(1 + \text{epoch}/5)$, shrinks the
cloud, and after 50 epochs the loss lies \num{4.4e-4} above the least.
The rate must fall slowly enough that the steps still add up to the
distance the weights must travel, and fast enough that the wandering
dies away; schedules of this kind, and the batch size, are among the
settings that every training run of planned Chapter~19 and after reports.

Each run of the figure evaluated the gradient of one reading
\num{100050} times; the full-batch run took 50 steps with them, the
mini-batch runs \num{10050}. When the readings are many and alike, as
the frames of a long run are, a mini-batch gradient is nearly as good as
the full one at a small fraction of the cost, and taking many cheap
noisy steps is faster than taking a few exact ones.

\begin{figure}[htb]
\centering
\includegraphics{ch18_learning/sgd}
\caption{Gradients from mini-batches. The force of the A--B bond at
\SI{300}{\kelvin} (2001 readings) fitted by a straight line in the
scaled stretch. (a) The mean squared miss over all the readings against
the epoch, for gradient descent on all of them, for batches of ten at a
fixed rate, and for batches of ten at a rate that falls, with the least
loss (dashed). (b) The gradient of the slope's weight from 1000 batches
of ten at the least loss, against a Gaussian of spread $\sigma_g/\sqrt{10}$
(dashed).}
\label{fig:lo-sgd}
\end{figure}

\begin{keyidea}
A mini-batch of $b$ readings estimates the gradient with an error that
falls as $\sigma_g/\sqrt b$, at $b/n$ of the cost. At a fixed rate the
weights end in a cloud about the minimum whose size the rate sets; a
rate that falls slowly shrinks it.
\end{keyidea}

\notebook{18}{change the batch size and the schedule and watch the loss
settle}

\begin{selfcheck}
Batches of 10 give a gradient that scatters by 0.18. What batch size
would bring the scatter to 0.06, and how many steps would an epoch of
2001 readings then take?
\end{selfcheck}
